\documentclass[12pt, a4paper]{article}

% ============================================================
% Packages
% ============================================================
\usepackage[utf8]{inputenc}
\usepackage[T1]{fontenc}
\usepackage{amsmath, amssymb}
\usepackage{graphicx}
\usepackage[margin=2.5cm]{geometry}
\usepackage{hyperref}
\usepackage{booktabs}
\usepackage{natbib}
\usepackage{setspace}
\usepackage{lipsum} % Remove this — only used for placeholder text

% ============================================================
% Metadata
% ============================================================
\title{Your Paper Title Here}
\author{
  First Author\textsuperscript{1} \and
  Second Author\textsuperscript{2}
}
\date{May 2026}

% ============================================================
% Document
% ============================================================
\begin{document}

\maketitle

\begin{abstract}
  Replace this with your abstract. This template is designed to be
  version-controlled with Git and automatically compiled to PDF via
  GitHub Actions. Every push to the \texttt{paper/} directory triggers
  a build, and every GitHub Release attaches the compiled PDF.
\end{abstract}

\tableofcontents
\newpage

% ============================================================
\section{Introduction}
\label{sec:introduction}
% ============================================================

Introduce your research question and motivation here.

This paper analyses historical temperature trends using data from the
Australian Bureau of Meteorology. Our approach builds on prior work by
\citet{example2024} and extends it to include seasonal decomposition.

% ============================================================
\section{Methods}
\label{sec:methods}
% ============================================================

Describe your methodology here.

\subsection{Data Collection}

We collected daily temperature records from 2000--2023 across three
Australian capital cities: Sydney, Melbourne, and Brisbane.

\subsection{Statistical Analysis}

We used linear regression to detect long-term trends:

\begin{equation}
  T(t) = \beta_0 + \beta_1 t + \epsilon
  \label{eq:trend}
\end{equation}

where $T(t)$ is the mean annual temperature, $t$ is the year, and
$\epsilon$ is the residual error term.

% ============================================================
\section{Results}
\label{sec:results}
% ============================================================

Present your results here. Reference figures and tables:

% Uncomment when you have a figure:
% \begin{figure}[htbp]
%   \centering
%   \includegraphics[width=0.8\textwidth]{../figures/temperature_trends.pdf}
%   \caption{Annual mean temperature trends for three Australian cities, 2000--2023.}
%   \label{fig:trends}
% \end{figure}

\begin{table}[htbp]
  \centering
  \caption{Temperature trend coefficients by city.}
  \label{tab:trends}
  \begin{tabular}{lcc}
    \toprule
    City      & Slope (\textdegree C/year) & $R^2$ \\
    \midrule
    Sydney    & 0.023 & 0.67 \\
    Melbourne & 0.019 & 0.58 \\
    Brisbane  & 0.031 & 0.72 \\
    \bottomrule
  \end{tabular}
\end{table}

% ============================================================
\section{Discussion}
\label{sec:discussion}
% ============================================================

Interpret your results here. Compare with prior work, discuss
limitations, and suggest future directions.

% ============================================================
\section{Conclusion}
\label{sec:conclusion}
% ============================================================

Summarise your key findings and their implications.

% ============================================================
% References
% ============================================================
\bibliographystyle{plainnat}
\bibliography{references}

% ============================================================
% Appendix (optional)
% ============================================================
% \appendix
% \section{Supplementary Material}

\end{document}
